% Zdroj: PDF priklady-podzim-2022.pdf, fyzická str. 7-8. Značka: specializace UI-SU. Zařazeno: S3. Číselné výsledky ověřeny numericky.
\section{Přeučení a regularizace}
\szzotazka{podzim 2022}{18}{Přeučení a regularizace}

\noindent\textbf{Zadání.}\par
Předpokládejme, že chceme trénovat lineární model pro regresi.

\begin{enumerate}
  \item Definujte střední čtvercovou chybu (MSE). Definujte upravenou chybovou funkci používanou při L2 regularizaci.
  \item Mějme dva lineární modely $M_1$, $M_2$ pro data se čtyřmi atributy:
    \[
      M_1(x) = x_1 - 2x_2 + 4x_3 - x_4 + 3, \qquad
      M_2(x) = x_2 - 5x_3 + x_4 - 1.
    \]
    Oba modely mají pro daná data stejnou MSE. Který má menší hodnotu upravené chybové funkce pro L2 regularizaci?
  \item Popište alespoň jeden další způsob pro zlepšení generalizačních schopností modelů strojového učení.
\end{enumerate}

\medskip
\noindent\textbf{Řešení.} \textit{(V~PDF bez náčrtu řešení; vlastní vzorové řešení, výpočet ověřen numericky.)}\par
\begin{enumerate}[nosep]
  \item Střední čtvercová chyba je
    \[
      \mathrm{MSE} = \frac{1}{n}\sum_{i=1}^{n} (y_i - \hat{y}_i)^2,
    \]
    kde $\hat{y}_i$ je predikce modelu pro $i$-tý vstup a~$y_i$ skutečná hodnota. Při L2 regularizaci se k~chybě přidá penalizace za velikost vah:
    \[
      \mathrm{MSE} + \lambda \sum_{j} w_j^2, \qquad \lambda > 0,
    \]
    kde $w_j$ jsou váhy modelu a~$\lambda$ je síla regularizace. Absolutní člen (bias / intercept) se obvykle \emph{ne}regularizuje.
  \item Regularizační člen závisí jen na vahách (bez biasu). Pro $M_1$ jsou váhy $(1,-2,4,-1)$, takže $\sum w_j^2 = 1+4+16+1 = 22$. Pro $M_2$ jsou váhy $(0,1,-5,1)$, takže $\sum w_j^2 = 0+1+25+1 = 27$. Při stejné MSE má proto \emph{menší} regularizovanou chybu model $\boldsymbol{M_1}$ ($22 < 27$).

    Poznámka: kdyby se regularizoval i~absolutní člen, vyšlo by $M_1$: $22 + 3^2 = 31$ a~$M_2$: $27 + (-1)^2 = 28$, a~vyhrál by $M_2$. Právě proto je podstatné, že se bias neregularizuje.
  \item Generalizaci lze zlepšit například získáním více trénovacích dat, výběrem či redukcí příznaků (méně parametrů), předčasným zastavením (early stopping), dropoutem u~neuronových sítí, kombinováním modelů (ensembling) nebo jinou regularizací (L1). Vhodné hyperparametry (např.\ $\lambda$) se volí křížovou validací.
\end{enumerate}

